\noindent340\hfill\textbf{REMARKS.}\hfill\mbox{}
\begin{center}\textbf{REMARKS.}\end{center}

\par\noindent\hspace*{\parindent}The preceding unfinished adjustment is the sole entry in a notebook entitled »Rechnungen in Beziehung auf die trigonometrischen Messungen«. It is noteworthy because the polygon equations for a ring system are set up here for the first time. The 3 polygon equations consist of the angle equation for the inner polygon of the ring and the two condition equations that ensure that a point's coordinates recover their original values upon returning to it along a traverse formed by sides of the ring system. Equation 2) on p.~337 expresses these two conditions and can be derived as follows.

\par\noindent\hspace*{\parindent}Calculating from point 1 along the traverse 1. 21. 3. 4. 5. 6. 9. 10. 11. 12. 15. 16. 18. 19. 20 gives rectangular coordinates for points 21 and 1. In general, before adjustment, the resulting coordinates differ from the starting values for 21 and 1 and may be taken to correspond to points \(21^*\) and \(1^*\). Let \(k\) have position \(z_k=x_k+\mathrm{i}y_k\). Then
\[
\begin{aligned}
z_{1^*}-z_1={}&(z_{21}-z_1)+(z_3-z_{21})+(z_4-z_3)+\ .\ .\ .\\
&+(z_{20}-z_{19})+(z_{21^*}-z_{20})+(z_{1^*}-z_{21^*}).
\end{aligned}
\]
\par\noindent\hspace*{\parindent}To bring \(z_{1^*}\) to \(z_1\), the sides and angles of the traverse are varied differentially, with point 1 held fixed.
\[
\delta z_{1^*}=\delta(z_{21}-z_1)+\delta(z_3-z_{21})+\ .\ .\ .+\delta(z_{1^*}-z_{21^*}).
\]
\par\noindent\hspace*{\parindent}If \(\mathrm{s}_{\lambda\mu}\) is the length and \(\tau_{\lambda.\mu}\) the azimuth of side \(\lambda.\mu\), then
\[
z_\mu-z_\lambda=\mathrm{s}_{\lambda.\mu}\mathrm{e}^{\mathrm{i}\tau_{\lambda.\mu}}
\]
\noindent and
\[
\delta(z_\mu-z_\lambda)=(z_\mu-z_\lambda)(\delta\log\mathrm{s}_{\lambda.\mu}+\mathrm{i}\,.\,\delta\tau_{\lambda.\mu}),
\]
\noindent so that
\[
1)\quad\delta z_{1^*}=\Sigma(z_\mu-z_\lambda)\delta\log\mathrm{s}_{\lambda.\mu}+\mathrm{i}\Sigma(z_\mu-z_\lambda)\delta\tau_{\lambda.\mu},
\]
\noindent where the sums run over the points of the traverse. Here and in what follows the logarithm is the hyperbolic logarithm.

\par\noindent\hspace*{\parindent}For the moment, denote the angle formed by the successive sides \(\lambda.\mu\) and \(\mu.\nu\) by \(w_\mu\). Then
\[
\tau_{\mu.\nu}=\tau_{\lambda.\mu}+w_\mu\pm180^0
\]
\noindent and hence
\[
\delta\tau_{\mu.\nu}=\delta\tau_{\lambda.\mu}+\delta w_\mu;
\]
\noindent Using this equation to express each \(\delta\tau\) successively in terms of \(\delta\tau_{1.21}\), we obtain for the traverse
\[
\begin{aligned}
\Sigma(z_\mu-z_\lambda)\delta\tau_{\lambda.\mu}={}&(z_{1^*}-z_1)\delta\tau_{1.21}+(z_{1^*}-z_{21})\delta w_{21}\\
&+(z_{1^*}-z_3)\delta w_3+\ .\ .\ .+(z_{1^*}-z_{20})\delta w_{20}\\
&+(z_{1^*}-z_{21^*})\delta w_{21^*}.
\end{aligned}
\]
\par\noindent\hspace*{\parindent}\(z_{1^*}-z_1\) may, in comparison with the other differences \(z_{1^*}-z_\lambda\), be regarded
\par\clearpage
\noindent\textbf{RING OF TRIANGLES AROUND OLDENBURG.}\hfill341\par\medskip
\enlargethispage{1.5\baselineskip}
\begingroup\small\setlength{\jot}{1pt}\setlength{\abovedisplayskip}{4pt}\setlength{\belowdisplayskip}{4pt}
\setlength{\abovedisplayshortskip}{2pt}\setlength{\belowdisplayshortskip}{3pt}
\noindent as a small quantity of the first order. The first term may therefore be neglected in comparison with the others, and \(z_{1^*}\) and \(z_{21^*}\) may be replaced in the coefficients by \(z_1\) and \(z_{21}\). The figure on p.~332 also gives
\[
\begin{gathered}
w_{21}=a^0+b^0,\quad w_3=360^0-c^0,\quad w_4=d^0,\\
w_5=360^0-e^0,\quad w_6=f^0+g^0+h^0,\ \text{and so on,}\\
w_{20}=360^0-u^0,\qquad w_{21}=v^0
\end{gathered}
\]
\noindent so that
\[
\begin{aligned}
&2)\quad\Sigma(z_\mu-z_\lambda)\delta\tau_{\lambda.\mu}=(z_1-z_{21})(\delta a^0+\delta b^0+\delta v^0)\\
&-(z_1-z_3)\delta c^0+(z_1-z_4)\delta d^0-(z_1-z_5)\delta e^0\\
&+(z_1-z_6)(\delta f^0+\delta g^0+\delta h^0)\!+\ .\ .\ .-(z_1-z_{20})\delta u^0.
\end{aligned}
\]
\par\noindent\hspace*{\parindent}The first sum in equation (1) can also be rearranged.

\par\noindent\hspace*{\parindent}From
\[
\begin{aligned}
\frac{\mathrm{s}_{21.3}}{\mathrm{s}_{1.21}}&=\frac{\sin b'}{\sin b''}\,.\,\frac{\sin a'}{\sin a''}\\[7pt]
\frac{\mathrm{s}_{3.4}}{\mathrm{s}_{21.3}}&=\frac{\sin c'}{\sin c''}\\[5pt]
&\text{and so on.}
\end{aligned}
\]
\noindent it follows that
\[
\begin{aligned}
\delta\log\mathrm{s}_{21.3}={}&\delta\log\mathrm{s}_{1.21}+\operatorname{cotang}a'\,.\,\delta a'-\operatorname{cotang}a''\,.\,\delta a''\\
&+\operatorname{cotang}b'\,.\,\delta b'-\operatorname{cotang}b''\,.\,\delta b''\\
\delta\log\mathrm{s}_{3.4}={}&\delta\log\mathrm{s}_{21.3}+\operatorname{cotang}c'\,.\,\delta c'-\operatorname{cotang}c''\,.\,\delta c''\\
&.\quad.\quad.\quad.\quad.\quad.\quad.\quad.\quad.\quad.\quad.\quad.\quad.\quad.\quad.\quad.\quad.\quad.\quad.\quad.\quad.\\
&.\quad.\quad.\quad.\quad.\quad.\quad.\quad.\quad.\quad.\quad.\quad.\quad.\quad.\quad.\quad.\quad.\quad.\quad.\quad.\quad.\\
\delta\log\mathrm{s}_{20.21^*}={}&\delta\log\mathrm{s}_{19.20}+\operatorname{cotang}u'\,.\,\delta u'-\operatorname{cotang}u''\,.\,\delta u''\\
\delta\log\mathrm{s}_{21^*.1^*}={}&\delta\log\mathrm{s}_{20.21^*}+\operatorname{cotang}v'\,.\,\delta v'-\operatorname{cotang}v''\,.\,\delta v''.
\end{aligned}
\]
\par\noindent\hspace*{\parindent}It follows that
\[
\begin{aligned}
&\Sigma(z_\mu-z_\lambda)\delta\log\mathrm{s}_{\lambda.\mu}=(z_{1^*}-z_1)\delta\log\mathrm{s}_{1.21}\\
&+(z_{1^*}-z_{21})\bigl(\operatorname{cotang}a'\,.\,\delta a'-\operatorname{cotang}a''\,.\,\delta a''\\
&\qquad+\operatorname{cotang}b'\,.\,\delta b'-\operatorname{cotang}b''\,.\,\delta b''\bigr)\\
&+(z_{1^*}-z_3)(\operatorname{cotang}c'\,.\,\delta c'-\operatorname{cotang}c''\,.\,\delta c'')+\ .\ .\ .\\
&+(z_{1^*}-z_{20})(\operatorname{cotang}u'\,.\,\delta u'-\operatorname{cotang}u''\,.\,\delta u'')\\
&+(z_{1^*}-z_{21^*})(\operatorname{cotang}v'\,.\,\delta v'-\operatorname{cotang}v''\,.\,\delta v''),
\end{aligned}
\]
\noindent or, again replacing \(z_{1^*}\) and \(z_{21^*}\) by \(z_1\) and \(z_{21}\) as before,
\[
\begin{aligned}
&3)\quad\Sigma(z_\mu-z_\lambda)\delta\log\mathrm{s}_{\lambda.\mu}=\\
&\quad(z_1-z_{21})\bigl(\operatorname{cotang}a'\,.\,\delta a'-\operatorname{cotang}a''\,.\,\delta a''\\
&\qquad+\operatorname{cotang}b'\,.\,\delta b'-\operatorname{cotang}b''\,.\,\delta b''\\
&\qquad+\operatorname{cotang}v'\,.\,\delta v'-\operatorname{cotang}v''\,.\,\delta v''\bigr)\\
&+(z_1-z_3)(\operatorname{cotang}c'\,.\,\delta c'-\operatorname{cotang}c''\,.\,\delta c'')\\
&+(z_1-z_4)(\operatorname{cotang}d'\,.\,\delta d'-\operatorname{cotang}d''\,.\,\delta d'')+\ .\ .\ .\\
&+(z_1-z_{20})(\operatorname{cotang}u'\,.\,\delta u'-\operatorname{cotang}u''\,.\,\delta u'').
\end{aligned}
\]
\par\noindent\hspace*{\parindent}Thus, when the angle corrections are given in seconds of arc, equation (1) yields
\[
4)\quad\GaussOldenburgLeftBrace{\begin{aligned}
&206265\,.\,\delta z_{1^*}=\\
&\quad(z_1-z_{21})\bigl\{\operatorname{cotang}v'\,.\,\delta v'-\operatorname{cotang}v''\,.\,\delta v''\\
&\qquad+\operatorname{cotang}a'\,.\,\delta a'-\operatorname{cotang}a''\,.\,\delta a''\\
&\qquad+\operatorname{cotang}b'\,.\,\delta b'-\operatorname{cotang}b''\,.\,\delta b''\\
&\qquad+\mathrm{i}(\delta v^0+\delta a^0+\delta b^0)\bigr\}\\
&+(z_1-z_3)\{\operatorname{cotang}c'\,.\,\delta c'-\operatorname{cotang}c''\,.\,\delta c''-\mathrm{i}\,.\,\delta c^0\}\\
&+(z_1-z_4)\{\operatorname{cotang}d'\,.\,\delta d'-\operatorname{cotang}d''\,.\,\delta d''+\mathrm{i}\,.\,\delta d^0\}\\
&.\quad.\quad.\quad.\quad.\quad.\quad.\quad.\quad.\quad.\quad.\quad.\quad.\quad.\quad.\quad.\quad.\quad.\\
&+(z_1-z_{20})\{\operatorname{cotang}u'\,.\,\delta u'-\operatorname{cotang}u''\,.\,\delta u''-\mathrm{i}\,.\,\delta u^0\}.
\end{aligned}}
\]
\par\noindent\hspace*{\parindent}For points 1 and \(1^*\) to coincide, we must set
\[
z_{1^*}+\delta z_{1^*}-z_1=0,\qquad \delta z_{1^*}=z_1-z_{1^*}=-\delta z_1
\]
\noindent if the prefixed \(\delta\) denotes a correction to be applied with positive sign.
\par\endgroup
\clearpage
\noindent342\hfill{\small\textbf{REMARKS.\quad RING OF TRIANGLES AROUND OLDENBURG.}}\par\medskip

\par\noindent\hspace*{\parindent}The side equation 1) on p.~336 is the condition that the value of side \(\mathrm{s}_{1.21}\) is recovered when the calculation is carried through the ring system. We have
\[
\frac{\mathrm{s}_{1^*.21^*}}{\mathrm{s}_{1.21}}=\frac{\sin a'}{\sin a''}\,.\,\frac{\sin b'}{\sin b''}\ .\ .\ .\frac{\sin v'}{\sin v''},
\]
\noindent hence
\[
\begin{aligned}
&\delta\log\mathrm{s}_{1^*.21^*}=-\delta\log\mathrm{s}_{1.21}=-\frac{\delta z_1-\delta z_{21}}{z_1-z_{21}}\\
&\quad=\operatorname{cotang}a'\,.\,\delta a'-\operatorname{cotang}a''\,.\,\delta a''+\ .\ .\ .\\
&\quad+\operatorname{cotang}v'\,.\,\delta v'-\operatorname{cotang}v''\,.\,\delta v''.
\end{aligned}
\]
\par\noindent\hspace*{\parindent}For Briggs's logarithm, and if the corrections \(\delta a',\ \delta a''\), etc. are to be obtained in seconds of arc, the left-hand side of this equation must be multiplied by \(\frac{206265}{\mathrm{Mod.}}\).

\par\noindent\hspace*{\parindent}In formulas 1) and 2) on pp.~336 and 337, the original has
\[
\frac{\delta(1-21)}{1-21}\ \text{and}\ \delta1\qquad\text{in place of}\quad-\frac{\delta(1-21)}{1-21}\ \text{and}\ -\delta1.
\]
\par\noindent\hspace*{\parindent}In 1834, according to the »Generalbericht über die mitteleuropäische Gradmessung für das Jahr 1865«, pp.~22/23, \textsc{Gauss} communicated to Oldenburg's survey director v.~\textsc{Schrenck} the 21 spheroidal triangles of the ring around Oldenburg, adjusted to \(180^0+\text{Excess}\). In a letter of 3~August 1835, \textsc{Gauss} wrote as follows, according to the same Generalbericht, pp.~23/24.

\par\noindent\hspace*{\parindent}»As for the relative accuracy of the triangles, you are mistaken, Sir, in contrasting triangles 2--8 [in the figure on p.~332 these are triangles \(q,\ p,\ o,\ n,\ m,\ l,\ k\)] with the remaining triangles 9--21 and 1 [in the figure \(i\) to \(a\), \(v\) to \(r\)]. The proper comparison is that 3--16 [in the figure \(p\) to \(b\)] are \textls[120]{much} less accurate than 17--21 [\(a,\ v,\ u,\ t,\ s\)], 1 [\(r\)] and 2 [\(q\)]. I measured these last 7 triangles myself with 12-inch theodolites, taking the greatest care, using heliotrope light for every target, and taking all possible care to ensure the stability of the observing stations, 3 of which were at ground level. By contrast, the 14 other triangles were measured for a different purpose, with a less powerful instrument (an 8-inch theodolite), in much less time, without heliotrope light and sometimes at very unfavourable stations, such as the towers of Twistringen and Asendorf and perhaps some of the other towers as well. Thus, since I consider the first 7 triangles to have been measured as precisely as the state of the art permits, I would assign the 14 remaining triangles only \(\frac{1}{3}\) of that accuracy, or only \(\frac{1}{9}\) of that weight.«

\par\noindent\hspace*{\parindent}If the sum of the angles of the heptagon 4. 6. 10. 12. 16. 19. 21 is formed from the angles communicated by \textsc{Gauss} to v.~\textsc{Schrenck}, then, according to the Generalbericht (p.~24), the closure error is \(14\overset{\prime\prime}{,}900\).
\begin{flushright}\textsc{Krüger.}\end{flushright}
\begin{center}\rule{.26\textwidth}{.4pt}\end{center}
